\section{Risiken und Herausforderungen}\label{sec:risiken}
Den Potenzialen stehen Befunde gegenüber, die dieselben Praktiken betreffen. Sie ordnen sich in zwei Gruppen: Qualität und Sicherheit des erzeugten Codes sowie Wirkungen auf Kompetenz, Prozess und Recht.

\subsection{Qualitäts- und Sicherheitsrisiken}\label{subsec:qualitaetsrisiken}
Der erzeugte Code ist nicht durchgängig sicher. Eine Untersuchung von 1.689 durch einen Assistenten erzeugten Programmen wies rund 40\,\% als verwundbar aus.\footcite[\vglf][\pagef 754]{Pearce.2022} Eine Messstudie über 200 Aufgaben, vier Sprachen und fünf Modellfamilien zeigt zudem Sprachabhängigkeit: In C und C++ treten Speicherfehler, fest eingebettete Geheimnisse und kryptografische Fehlanwendungen häufiger auf, moderne Sicherheitsmerkmale bleiben oft ungenutzt.\footcite[\vglf][\pagef 7300]{Kharma.2026} In einer Nutzerstudie mit 47 Teilnehmern schrieb die Gruppe mit Assistenz signifikant unsichereren Code und hielt ihn zugleich häufiger für sicher.\footcite[\vglf][\pagef 2785]{Perry.2023} Solches Übernehmen fehlerhafter Modellausgaben wird als \gls{overreliance} bezeichnet.\footcite[\vglf][\pagef 4]{Shen.2026} Hinzu tritt der Korrekturaufwand für Vorschläge, die fast richtig sind: 66\,\% der Befragten einer Entwicklerumfrage nennen ihn als Hauptärgernis, 45\,\% das \glslink{debugging}{Debuggen} KI-erzeugten Codes.\footcite[\vglf][\pagef 71]{Swaraj.2026} Auf Prozessebene sinken Auslieferungsstabilität ($-$7,2\,\%) und Durchsatz ($-$1,5\,\%) je 25\,\% höherer KI-Adoption; als Erklärung vermutet der Report größere Änderungspakete und erinnert an das Prinzip kleiner Losgrößen.\footcite[\vglf][\pagef 39 f.]{DORA.2024} Daraus ergibt sich eine Rückkopplung: Code Review und Definition of Done sollen solche Fehler abfangen, setzen aber die prüfende Aufmerksamkeit voraus, die Overreliance gerade schwächt.

\subsection{Kompetenz-, prozess- und rechtsbezogene Risiken}\label{subsec:kompetenzrisiken}
Die zweite Gruppe betrifft Menschen und Prozesse. Ein randomisiertes Experiment zum Erlernen einer unbekannten Programmbibliothek -- durchgeführt von einem Modellanbieter und bislang als Vorabfassung veröffentlicht -- findet, dass KI-Unterstützung Verständnis, Code-Lesen und Debugging beeinträchtigt, ohne im Mittel Zeit zu sparen; wer die Aufgabe vollständig delegierte, gewann Produktivität auf Kosten des Lernens.\footcite[\vglf][\pagef 1]{Shen.2026} Dazu passt, dass rund 20\,\% der befragten Entwickler geringeres Vertrauen in die eigenen Fähigkeiten berichten und nur 17\,\% eine bessere Zusammenarbeit im Team beobachten.\footcite[\vglf][\pagef 71]{Swaraj.2026} Daraus folgt ein Risiko für die gemeinsame Verantwortlichkeit, eine der zwölf Vorgehensweisen des \gls{xp}:\footcite[\vglf][\pagef 152 f.]{Preussig.2020} Sie setzt voraus, dass Teammitglieder fremden Code lesen und ändern können -- genau die Fähigkeit, die nach diesen Befunden leidet. Hinzu kommt eine Wahrnehmungslücke: Dieselben Entwickler, die 19\,\% länger brauchten, schätzten sich im Nachhinein um 20\,\% schneller ein; Aufwandsschätzungen und Velocity-Werte dürften dadurch verzerrt werden.\footcite[\vglf][\pagef 2]{Becker.2025} Rechtlich und organisatorisch bleibt die Vertraulichkeit offen. Fehlen verbindliche Richtlinien, entscheiden Beschäftigte nach eigenem Ermessen; Code mit Schlüsseln oder personenbezogenen Daten darf Werkzeugen ohne gesicherte Datenhoheit nicht übergeben werden.\footcite[\vglf][\pagef 298]{Neumann.2026}
